% AI-produced by GPT-6 Astra (OpenAI, ChatGPT Pro mode), 19 November 2026.
% A method obstruction only. No movement credit on the original A.5 gate.
\subsection{A product bundle need not admit a Thom stratification}

The remaining maximal-locus question cannot be replaced by recognition of
an ambient Thom stratification. We prove this even for a proper polynomial
map with smooth algebraic source and closed manifold fibres. We use the
tangent-kernel convention of
\cite[Definitions 2.1--2.2]{A5-GimenezConejeroLeNuno2023}: each source
stratum maps submersively to a target stratum, and for every incident
source pair the limiting tangent kernels contain the kernel on the
lower stratum. Arbitrary finite compatible refinements of both source
and target are allowed.

\begin{saresult}{A5.015}{Proved}{A smooth-source product bundle with no finite ambient Thom stratification}
Let $R$ be any real closed field, equipped with any o-minimal expansion.
Put
\[
 M=S^2_R=\{x=(x_1,x_2,x_3)\in R^3:
                    x_1^2+x_2^2+x_3^2=1\},\qquad h(x)=x_1,
\]
and let $X=R^2\times M\times R$, with coordinates $(a,b,x,u)$.
The polynomial map
\[
 f:X\longrightarrow R^3,\qquad
 f(a,b,x,u)=(a,b,u^3-bh(x))
 \tag{A5.15.1}
\]
is definably proper and globally semialgebraically trivial, with closed
fibre $S^2_R$. Its source is a smooth algebraic $5$-manifold. Nevertheless,
for every $a_0\in R$, no open neighbourhood of $(a_0,0,0)$ admits a
finite definable ambient $C^1$ Thom $a_f$ stratification for $f$.
This is a method obstruction, not a counterexample to A.5: the maximal
open manifold-and-bundle locus of this map is all of $R^3$.
\end{saresult}
\begin{proof}
We give the product map, the properness bound, and the argument excluding
every finite stratification separately. Throughout, compactness means
definable compactness; the argument does not use ordinary compactness
over a non-Archimedean field.

\medskip\noindent\textbf{1. A global product and properness.}
The map
\[
 H:R^3\times M\longrightarrow X,\qquad
 H(a,b,c,x)=(a,b,x,\sqrt[3]{c+bh(x)})
 \tag{A5.15.2}
\]
is a semialgebraic homeomorphism. The real cube root exists uniquely in
every real closed field and is continuous and semialgebraic. The inverse
is
\[
 (a,b,x,u)\longmapsto(a,b,u^3-bh(x),x),
\]
and $fH=\operatorname{pr}_{R^3}$. Thus all fibres are semialgebraically
homeomorphic to $S^2_R$, and the claimed bundle locus is the whole base.

For completeness, if $K\subset R^3$ is closed, bounded and definable,
choose $C\ge1$ such that $|a|,|b|,|c|\le C$ on $K$.
At every point of $f^{-1}(K)$ we have
\[
 |u|^3=|c+bh(x)|\le |c|+|b|\le2C,
\]
since $|h|\le1$ on $M$. Hence $|u|\le2C$, and the other coordinates
are bounded as well. The set $f^{-1}(K)$ is closed, because $X$ is a
closed algebraic subset of its Euclidean ambient space and $f$ is
continuous. It is therefore closed and bounded, proving definable
properness. Over $R=\mathbb R$ this is also ordinary properness.
The source is smooth because the defining sphere equation has nonzero
gradient at every point of $M$.

\medskip\noindent\textbf{2. A central tangent direction and an approaching kernel.}
Fix $a_0\in R$ and suppose that a finite definable $C^1$ Thom
stratification exists over an open neighbourhood $U$ of $(a_0,0,0)$.
We do not assume in advance that the fibre
\[
 E=\{(a_0,0,x,0):x\in M\}
\]
is a union of source strata. Its open subset where $|x_1|<1$ has
dimension two. Finiteness and definable dimension imply that some
source stratum $T$ has intersection with this subset containing a
nonempty relatively open subset of $E$. Choose
\[
 q=(a_0,0,x^0,0)\in T,\qquad |x^0_1|<1,
\]
in such a relative interior. A neighbourhood of $q$ in $E$ lies in
$T$, so the central sphere tangent plane satisfies
\[
 C_0=\{(0,0,v,0):v\in T_{x^0}M\}
                   \subseteq\ker d(f|_T)_q.
 \tag{A5.15.3}
\]
For $t>0$ consider the semialgebraic curve
\[
 q_t=(a_0,t,x^0,0),\qquad
 f(q_t)=(a_0,t,-t h(x^0)).
 \tag{A5.15.4}
\]
As $t\to0^+$, these points tend to $q$ and their target values tend
to $(a_0,0,0)$; hence a tail lies in $f^{-1}(U)$.

A tangent vector to $X$ has coordinates $(\alpha,\beta,v,\eta)$
with $v\in T_xM$. Direct differentiation gives
\[
 df_{(a,b,x,u)}(\alpha,\beta,v,\eta)
   =\bigl(\alpha,\beta,
           3u^2\eta-h(x)\beta-b(dh)_xv\bigr).
 \tag{A5.15.5}
\]
Consequently the full vertical tangent plane at every $q_t$ is
the same subspace
\[
 V=\ker df_{q_t}
   =\{(0,0,v,\eta):v\in T_{x^0}M,
                         (dh)_{x^0}v=0,\ \eta\in R\}.
 \tag{A5.15.6}
\]
The differential $(dh)_{x^0}$ is nonzero, since $|x^0_1|<1$.
Thus $V$ has dimension two, and $df_{q_t}$ has rank three.
More explicitly, the vector
\[
 v_0=(1-(x^0_1)^2,-x^0_1x^0_2,-x^0_1x^0_3)
 \tag{A5.15.7}
\]
is tangent to $M$ at $x^0$: its scalar product with $x^0$ is
$x^0_1(1-\lVert x^0\rVert^2)=0$. It has
$(dh)_{x^0}v_0=1-(x^0_1)^2>0$. Therefore
\[
 e=(0,0,v_0,0)\in C_0,\qquad e\notin V.
 \tag{A5.15.8}
\]

\medskip\noindent\textbf{3. Excluding every finite compatible refinement.}
By o-minimality, a tail of $q_t$ lies in one source stratum $S$.
Let $K_t=\ker d(f|_S)_{q_t}$. Since $T_{q_t}S\subset T_{q_t}X$,
we have $K_t\subseteq V$. The dimension of $K_t$ is constant:
$f|_S$ is a submersion to one fixed target stratum.

These subspaces have a definable limit $L$ as $t\to0^+$. Here is
an argument valid over the given real closed field. Represent $K_t$
by its orthogonal projection matrix $P_t$ in the fixed ambient
coordinate space. The matrix is definable and its entries are bounded.
Every bounded definable function of one variable has a finite limit
at an endpoint, so $P=\lim_{t\to0^+}P_t$ exists. Taking limits in
$P_t^2=P_t=P_t^{\mathsf T}$ shows that $P$ is again an orthogonal
projection. Its trace is the same fixed integer $\dim K_t$, so its
image $L$ has that dimension. Since each image of $P_t$ is contained
in the fixed space $V$, the image of $P$ is contained in $V$ as well.
This gives
\[
 L=\lim_{t\to0^+}K_t\subseteq V.
 \tag{A5.15.9}
\]
No assertion of sequential compactness has been used.

If $S=T$, the kernel bundle of the constant-rank $C^1$ map
$f|_T$ is continuous at $q$. Hence $L=\ker d(f|_T)_q$,
contradicting (A5.15.3) and (A5.15.8), because $L\subseteq V$.
If $S\ne T$, then $q\in T\cap\overline S$ and the frontier
condition gives $T\subseteq\overline S$. Thom's condition for this
incident pair would require
\[
 Re\subseteq\ker d(f|_T)_q\subseteq L\subseteq V,
\]
again contradicting $e\notin V$. This rules out every finite
compatible definable source and target stratification on the
specified neighbourhood, proving the result.
\end{proof}

The argument also works over $R^2$ after deleting the passive coordinate
$a$, and with $S^2$ replaced by a closed definable $C^1$ manifold carrying
a definable $C^1$ function with nonzero differential somewhere. We keep
the displayed $R^3$--$S^2$ formulation because these are exactly the first
dimensions not covered by A5.013--A5.014. The map is already an actual
global product in the definable category; its failure of a Thom
stratification therefore disproves the proposed necessity criterion,
while proving neither the truth nor the falsity of the original
maximal-locus assertion.

Our earlier rational graph construction and its full interval and sphere
proofs remain in
\path{agents/maximal_locus/notes/product_without_thom.tex} as route
history. They require no correction, but the polynomial construction
above has the stronger smooth-source property and a simpler proof.
The finite exact-algebra certificate is
\path{agents/maximal_locus/code/check_smooth_product_without_thom.py},
with full output in
\path{agents/maximal_locus/logs/check_smooth_product_without_thom.json}.
It checks the displayed differential, approach curve, sphere tangent
vector and rank minor. The proof, rather than a finite enumeration,
handles the quantifier over all stratifications. We claim no priority.



% AI-produced by GPT-6 Astra (OpenAI, ChatGPT Pro mode), 19 November 2026.
% Method/special-model result only. No movement credit on original A.5.
\subsection{Coherent approximation for one real-definable cell-like surface map}

We prove a finite height-leveling lemma, give the elementary dimension-three
link argument needed below, and apply them to an individual definable cell-like
map between closed surfaces over the real numbers. The first lemma works over
every real closed field. The last theorem uses classical topology over
$\mathbb R$ and is not asserted over arbitrary real closed fields or uniformly
in additional parameters. None of these claims closes the remaining A.5 gate.

\paragraph{Original-source inputs and precise locators.}
Rourke--Sanderson~\cite[Lemma 4.23, p.~58]{RourkeSanderson1972} levels a PL collar in a product by matching derived
subdivisions whose new vertices lie at a common positive height. We use that
finite construction for two possibly different nonnegative PL height functions
and supply the details. The same original gives the compact PL collar theorem
as Theorem 2.25, p.~24; the closed PL surface classification is used explicitly
in the remarks following Theorem 5.5, p.~63; the dimension-two PL Schoenflies
conclusion is recorded in the discussion of Problem 3.37, p.~47. Full original:
\url{https://webhomes.maths.ed.ac.uk/~v1ranick/papers/rsplx.pdf}.

Shiota~\cite[Theorem 2.2(1), p.~6]{A5-Shiota} gives exact definable
triangulation of a scalar function, with no reparametrization of its target.
Section 2 explicitly works over a real closed field. Full original:
\url{https://arxiv.org/pdf/1002.1508}.

Siebenmann~\cite[Approximation Theorem A, p.~271]{Siebenmann1972Approximation} applies to closed surfaces by condition (a).
Its Complement, p.~272, gives a continuous family of homeomorphisms with the
cell-like map as its endpoint. Full original:
\url{https://webhomes.maths.ed.ac.uk/~v1ranick/papers/siebappr.pdf}.
Both printed pp.~271--272 were visually reread from the original scan during
this audit. We use the continuous endpoint family only to establish that the
mapping cylinder is a topological three-manifold; the family is not claimed
to be definable.

\paragraph{Lemma 1 (finite leveling of a collared zero set).}
Let $R$ be a real closed field, let $Y$ be a compact polyhedron, and let
$g:Y\to[0,M]$ be PL. Set $L=g^{-1}(0)$ and suppose that $L$ admits a PL collar
in $Y$. Then there are $\epsilon>0$ and a PL homeomorphism
\[
 C:L\times[0,\epsilon]\longrightarrow g^{-1}([0,\epsilon])
 \quad\text{such that}\quad
 C(z,0)=z,\qquad g(C(z,s))=s.
 \tag{A5.16.1}
\]
In particular the given scalar height, rather than an unspecified collar
coordinate, is preserved.

\paragraph{Proof.}
Choose a PL collar $c:L\times[0,1]\to Y$, with $c(z,0)=z$, and write
$Y'=c(L\times[0,1])$. The image of $L\times[0,1)$ is an open neighborhood
of $L$ in $Y$. Its complement is a compact polyhedron disjoint from $L$,
so $g$ has a positive lower bound there. Thus a sufficiently small positive
$\epsilon$ has $g^{-1}([0,\epsilon])\subset Y'$.

Triangulate the finite PL diagram consisting of $c$, the domain height
$t:L\times[0,1]\to[0,1]$, and $g|Y'$. This gives finite complexes $K,J$
and a simplicial isomorphism
\[
 c:K\longrightarrow J,
 \qquad |K|=L\times[0,1],\quad |J|=Y',
\]
such that $t$ is linear on every simplex of $K$ and $g$ is linear on every
simplex of $J$. The zero vertices of $t$ correspond exactly to the zero
vertices of $g$, because both zero sets correspond to $L$ under $c$.
Shrink $\epsilon$ further so that $0<\epsilon<1$ and $\epsilon$ is less
than every positive number in the finite lists
\[
 \{t(v):v\in K^{(0)},\ t(v)>0\},\qquad
 \{g(c(v)):v\in K^{(0)},\ t(v)>0\}.
 \tag{A5.16.2}
\]

Call an original simplex of $K$ mixed if it has both a zero vertex and a
positive vertex. For each mixed simplex $\sigma$, its relative interior
meets $t^{-1}(\epsilon)$: positive barycentric coefficients may be chosen
whose weighted height is $\epsilon$, since the minimum vertex height is
zero and the maximum is greater than $\epsilon$. Similarly the relative
interior of $c(\sigma)$ meets $g^{-1}(\epsilon)$.

In decreasing order of dimension, stellarly subdivide every original mixed
simplex $\sigma$ at a point
\[
 a_\sigma\in\operatorname{relint}(\sigma)\cap t^{-1}(\epsilon),
\]
and subdivide $c(\sigma)$ at an independently chosen point
\[
 b_\sigma\in\operatorname{relint}(c(\sigma))\cap g^{-1}(\epsilon).
\]
Each original face is still present when it is to be starred: previous stars
were in relative interiors of larger faces or in other faces of the same
dimension. Paired stellar subdivisions therefore produce combinatorially
isomorphic complexes $K'$ and $J'$. Their canonical simplicial isomorphism
$c':K'\to J'$ sends each old vertex $v$ to $c(v)$ and each $a_\sigma$ to
$b_\sigma$, and fixes $L$ pointwise.

Every new vertex has height $\epsilon$. Every old positive vertex has height
greater than $\epsilon$ in the respective complex. No simplex of $K'$ can
contain both an old zero vertex and an old positive vertex. Indeed, these two
vertices would span an edge of $K'$; stellar subdivisions introduce no new
edges between old vertices, while every original zero--positive edge was
itself starred. The identical assertion holds in $J'$.

It follows that the sublevel set $t^{-1}([0,\epsilon])$ is exactly the
subcomplex of $K'$ with vertex heights $0$ or $\epsilon$. To see this
directly, a simplex containing a zero vertex contains no old positive vertex,
so all its heights are at most $\epsilon$. A simplex containing an old
positive vertex has no zero vertex, and its intersection with the sublevel
set is precisely the face spanned by its height-$\epsilon$ vertices.
The same description holds for $g^{-1}([0,\epsilon])$ in $J'$.

The isomorphism $c'$ identifies these two subcomplexes. On their vertices it
preserves height, since the only heights are $0$ and $\epsilon$; linearity
on each simplex then gives $g\circ c'=t$ throughout the band. Restricting
$c'$ gives (A5.16.1). All operations are finite PL operations, so the argument
works over $R$ and does not require a mesh whose size grows as the height
tends to zero. \hfill$\square$

\paragraph{Lemma 2 (the finite link argument in dimension three).}
Let $K$ be a finite simplicial complex over $\mathbb R$ whose underlying
space is a topological three-manifold, possibly with boundary. Assume that
the topological boundary is the underlying space of a subcomplex $B$.
Then $K$ is a PL three-manifold with boundary $B$.

\paragraph{Proof.}
First $K$ is pure of dimension three. An interior point of a maximal simplex
has a neighborhood of the dimension of that simplex; manifold dimension
and the existence of interior points in every manifold neighborhood force
every maximal simplex to have dimension three.

For a $k$-simplex $\sigma$ and a point $x\in\operatorname{relint}(\sigma)$,
the standard local homology formula is
\[
 H_i(|K|,|K|\setminus\{x\};\mathbb Z)
 \cong
 \widetilde H_{i-k-1}(\operatorname{lk}_K\sigma;\mathbb Z).
 \tag{A5.16.3}
\]
At an interior manifold point the left side is $\mathbb Z$ for $i=3$ and
zero otherwise. At a boundary point it is zero in every degree. Since $B$
is a subcomplex, each open simplex consists entirely of interior points or
entirely of boundary points.

If $\sigma$ is a two-simplex, its link is a nonempty finite set. Equation
(A5.16.3) says that the link has two points in the interior case and one point
in the boundary case. Thus every triangle has respectively two or one
incident tetrahedra.

If $\sigma$ is an edge, its link is a finite pure graph. Each vertex of this
graph corresponds to a triangle containing $\sigma$, so its degree is two
or one by the preceding paragraph. Equation (A5.16.3) makes the graph connected;
its first homology is $\mathbb Z$ in the interior case and zero in the
boundary case. A connected finite graph with all degrees one or two is
therefore a circle or an interval respectively. These are PL circles and
PL intervals, with their given finite triangulations.

For a vertex $v$, the link $S=\operatorname{lk}_K(v)$ is a compact pure
two-complex. The link of any vertex of $S$ is the link in $K$ of the
corresponding edge through $v$, hence is a circle or an interval. Therefore
$S$ is a compact PL surface, possibly with boundary. Equation (A5.16.3) makes
$S$ connected. If $v$ is interior, it has the integral homology of $S^2$.
The classification of compact PL surfaces makes it a PL two-sphere.

If $v$ is a boundary vertex, $S$ is acyclic. An acyclic compact connected
surface cannot be closed: an orientable closed surface has nonzero $H_2$,
and a nonorientable closed surface has nonzero integral $H_1$. Thus
$\partial S\ne\varnothing$. To include the boundary count explicitly, an
orientable compact connected surface of genus $r$ with $b>0$ boundary
components has
\[
 H_1(S;\mathbb Z)\cong\mathbb Z^{2r+b-1},
\]
while a nonorientable surface with $r\ge1$ crosscaps and $b>0$ boundary
components has $H_1(S;\mathbb Z)\cong\mathbb Z^{r+b-1}$.
Vanishing of $H_1$ forces the orientable case $r=0,b=1$; the nonorientable
case is impossible. Hence $S$ is a PL disk, and its boundary is connected.

For completeness this disk conclusion can be obtained using only the
closed-surface classification quoted above. Double $S$ along its $b$
boundary circles. Mayer--Vietoris gives $H_2(DS)\cong\mathbb Z^b$ and
$H_1(DS)\cong\mathbb Z^{b-1}$ because $S$ is acyclic. The double is a
connected closed PL surface, so $b=1$ and $DS$ is a PL sphere. The embedded
boundary circle is locally flat, and the two-dimensional PL Schoenflies
theorem makes either complementary closure, including $S$, a PL disk.

Thus every vertex link is the requisite PL sphere or PL ball. This is the
combinatorial link criterion for $K$ to be a PL three-manifold. The boundary
triangles are exactly those with one incident tetrahedron, so the PL boundary
is exactly the given topological boundary $B$. No assertion that an arbitrary
topological homeomorphism is PL is used. \hfill$\square$

\begin{saresult}{A5.016}{Proved}{Coherent definable approximation of one cell-like closed-surface map}
Fix an o-minimal expansion of $\mathbb R$. Let $F\subset\mathbb R^m$ and
$G\subset\mathbb R^n$ be compact definable closed topological surfaces,
and let $q:F\to G$ be a continuous definable surjection whose point inverses
are cell-like in the classical topological sense. Then there are
$\epsilon>0$ and a jointly definable continuous map
\[
 k:(0,\epsilon]\times G\longrightarrow F
\]
such that each $k_s$ is a homeomorphism and
\[
 \lim_{s\downarrow0}\sup_{y\in G}\|q(k_s(y))-y\|=0.
 \tag{A5.16.4}
\]
Equivalently, $q$ is the endpoint of a definable continuous family
$h_s:F\to G$ with $h_0=q$ and $h_s$ a homeomorphism for every $s>0$.
Consequently the associated double mapping-cylinder projection is definably
locally trivial at its central parameter.

\end{saresult}
\begin{proof}
\textbf{First establish the topological three-manifold.}
If $F=G=\varnothing$, the unique empty homeomorphism gives the conclusion,
with uniform error defined to be zero. If exactly one surface is empty,
the assumed surjection does not exist. We therefore assume both nonempty.
The assertion may be checked on components. Cell-like fibres are connected,
so components of $F$ map bijectively to components of $G$; both have finitely
many components. Assume the surfaces are connected.

By Siebenmann's Theorem A and its Complement there is a continuous map
$(t,x)\mapsto a_t(x)$ on $[0,1]\times F$ such that $a_1=q$ and
$a_t:F\to G$ is a homeomorphism for $0\le t<1$. Its continuity on a compact
domain gives uniform convergence $a_t\to q$ as $t\uparrow1$. For $s>0$ put
$b_s=a_{1-s}^{-1}$. The inverse family is jointly continuous for positive
$s$, as follows by taking inverses of the compact-domain homeomorphism
$(s,x)\mapsto(s,a_{1-s}(x))$ on any closed positive $s$-interval. Moreover
\[
 \sup_{y\in G}\|q(b_s(y))-y\|
 =\sup_{x\in F}\|q(x)-a_{1-s}(x)\|\longrightarrow0.
 \tag{A5.16.5}
\]

Consider the compact definable mapping cylinder
\[
 C_q=\{(s,sx,q(x)):0<s\le1,\ x\in F\}
       \cup\{(0,0,y):y\in G\},\qquad \tau(s,z,y)=s.
 \tag{A5.16.6}
\]
The map from $[0,1]\times G$ to $C_q$ given for $s>0$ by
\[
 (s,y)\longmapsto(s,sb_s(y),q(b_s(y)))
\]
and at $s=0$ by $(0,y)\mapsto(0,0,y)$ is continuous by boundedness of
$F$ and (A5.16.5), and is bijective. Compactness and the Hausdorff property
make it a homeomorphism. Thus $C_q$ is a topological three-manifold whose
boundary is the disjoint union of the levels $\tau=0$ and $\tau=1$.
This use of a classical family has established only this topological fact.

\medskip\noindent\textbf{Replace the height collar by a finite definable one.}
Apply Shiota's exact scalar triangulation theorem to
$\tau:C_q\to[0,1]$. There are a finite compact polyhedron $P$ and a
definable homeomorphism $\pi:P\to C_q$ such that $\ell=\tau\circ\pi$ is
PL. Subdivide so that $\ell$ is linear on each simplex. Since its values
lie in $[0,1]$, both levels $L_0=\ell^{-1}(0)$ and $L_1=\ell^{-1}(1)$
are subcomplexes. They triangulate precisely the topological boundary.
Lemma 2 makes $P$ a PL three-manifold, with boundary $L_0\sqcup L_1$.

The compact PL boundary component $L_0$ has a PL collar. Apply Lemma 1 to
$\ell$ to obtain a PL height-preserving collar
\[
 C:L_0\times[0,\epsilon]\longrightarrow\ell^{-1}([0,\epsilon]).
\]
Use the definable identification $j:G\to L_0$ determined by
$\pi(j(y))=(0,0,y)$. Then
\[
 \Phi(s,y)=\pi(C(j(y),s))
\]
is a definable height-preserving homeomorphism from $[0,\epsilon]\times G$
onto $\tau^{-1}([0,\epsilon])$, with the identity on the zero fibre.
For $s>0$ write
\[
 \Phi(s,y)=(s,s k_s(y),q(k_s(y))).
 \tag{A5.16.7}
\]
Division of the middle coordinate by $s$ makes $k$ jointly definable and
continuous for $s>0$; each $k_s$ is a homeomorphism because $\Phi$ identifies
the height-$s$ fibres. Uniform continuity of the last coordinate on the
compact domain $[0,\epsilon]\times G$ gives (A5.16.4).

Set $h_s=k_s^{-1}$ for $s>0$ and $h_0=q$. The inverse family is jointly
definable and continuous for positive $s$, and
\[
 \sup_{x\in F}\|h_s(x)-q(x)\|
 =\sup_{y\in G}\|y-q(k_s(y))\|\longrightarrow0.
\]
Hence the extension at $s=0$ is jointly continuous. Reflecting the positive
height collar using $|s|$ supplies the claimed double-cylinder chart.
\end{proof}

\paragraph{What a parameter-preserving version would additionally need.}
The finite proof of Lemma 1 does not automatically preserve an extra map
$p:Y\to A$, even if a PL collar $c$ already preserves $p$. In a mixed edge,
the domain point at height $\epsilon$ and the image point at height
$\epsilon$ can require different barycentric fractions. If $p$ varies
along that edge, those choices need not have the same $p$-coordinate.

There is a direct counterexample to the unrestricted parameter-preserving
claim. On $Y=[-1,1]\times[0,1]$ put $p(a,t)=a$, $u=|a|$, and
\[
 g(a,t)=
 \begin{cases}
 3t,&0\le t\le u/2,\\
 2u-t,&u/2\le t\le u,\\
 t,&u\le t\le1.
 \end{cases}
 \tag{A5.16.8}
\]
The function is continuous and PL, positive exactly when $t>0$. The zero
set has the standard product PL collar over $a$. Each fixed $a$ admits
a height-preserving collar near its zero point. Nevertheless no uniform
collar near $(a,g)=(0,0)$ can preserve both $p$ and $g$: for every sufficiently
small $a>0$ and every $a<s<3a/2$, the fibre over $(a,s)$ has the three
distinct points
\[
 t=s/3,\qquad t=2a-s,\qquad t=s,
\]
whereas the central fibres have one point.

This counterexample is not a counterexample to the family mapping-cylinder
problem. In that problem every positive-height fibre is already the same
surface $F$; (A5.16.8) violates that condition. It shows precisely why a
parameter-preserving collar of the zero family, or a fibrewise PL manifold
structure alone, cannot justify parameter-preserving height leveling.

For a family $q_a:F\to G$, the needed conclusion remains a jointly definable
family $k(a,s,y)$ with $q_a k(a,s,y)\to y$ jointly at $s=0$. To use a finite
PL argument, one would need a parameter-compatible triangulation and a
leveling construction whose stellar choices preserve the parameter map, or
another genuine family theorem exploiting the special positive-height
product. The present argument does not supply either assertion.

\paragraph{Status.}
Lemma 1 closes the previously isolated prescribed-height collar substep.
Lemma 2 avoids an unjustified appeal to uniqueness of PL structures applied
to an arbitrary topological homeomorphism. A5.016 treats one cell-like
surface map over $\mathbb R$. It is a special-model result, not a reduction
or closure of the whole A.5 problem; no original-gate movement credit is
claimed. The uniform family and arbitrary-real-closed-field conclusions
remain unproved here.

\paragraph{Finite exact certificate.}
The accompanying checker
\path{check_surface_collar_parameter_fold.py}
verifies all three affine pieces in (A5.16.8), their two boundary matching identities, and the three distinct preimages with their region inequalities. It records 21 exact checks passing in 0.06886558799305931 seconds. These checks certify the displayed parameter-fold algebra only; the finite PL collar, link argument, and cell-like surface theorem are proved above.

% AI-produced by GPT-6 Astra (OpenAI, ChatGPT Pro mode), 19 November 2026.
% Root and higher_comparison2 full-proof audits PASS. No original A5 movement.
\subsection{Uniform formats for surface approximants in a definable family}

Fix an o-minimal expansion $\mathcal R$ of the real field. Let $A$ be a
definable parameter set, let $F,G$ be compact definable closed surfaces,
and let $q:A\times F\to G$ be definable. Assume each $q_a$ is a continuous
cell-like surjection. Joint continuity in $a$ is not required for the
format lemma, although it will be required for the intended sewing problem.

If $A=\varnothing$, the asserted family is empty. If $F=G=\varnothing$,
use the unique empty homeomorphism for every parameter and take its uniform
error to be zero. If exactly one of $F,G$ is empty and $A$ is nonempty, no
surjection $q_a:F\to G$ as in the hypotheses exists. Thus the proof below
may assume $A,F,G$ nonempty.

\begin{saresult}{A5.017}{Proved}{Uniform definable formats for approximating a family of cell-like surface maps}
There is a single definable map
\[
 \psi:A\times(0,1]\times F\longrightarrow G
 \tag{A5.17.1}
\]
such that, for every $(a,\eta)$, the slice $\psi_{a,\eta}:F\to G$ is a
homeomorphism and
\[
 \sup_{x\in F}\|\psi_{a,\eta}(x)-q_a(x)\|<\eta.
 \tag{A5.17.2}
\]
The map $\psi$ need not be jointly continuous in $a$ or $\eta$. Its graph
is nevertheless defined by one formula with finitely many fixed parameters.
The claim is consequently a genuine finite-format existence assertion,
stronger than independent existence of an approximant for each $(a,\eta)$.

\end{saresult}
\begin{proof}
\textbf{1. First-order scalar-triangulation certificates.}
For each $a$ form the compact definable cylinder $C_a=C_{q_a}$ and its scalar
height $\tau_a:C_a\to[0,1]$ as in A5.016. The predicates that the slices
$q_a$ are continuous surjections and the sets $C_a$ are closed and bounded
are first-order. They therefore persist in elementary extensions of
$\mathcal R$. No cell-likeness assertion is used at this step.

Consider all possible finite certificates for exact scalar triangulations
of these heights. A certificate specifies:
\begin{enumerate}
\item a finite abstract simplicial complex $K$ and coordinates for a geometric
realization $P$ of $K$;
\item a fixed formula $\varphi(u,z;b)$, with a finite tuple $b$ of freely
quantified coefficients, defining the graph of a homeomorphism
$\pi:P\to C_a$;
\item vertex heights, with the assertion that $\tau_a\circ\pi$ is their
affine interpolation on each simplex.
\end{enumerate}
Here is an explicit certificate predicate. Let $K$ have vertices
$1,\ldots,N$ and let $v=(v_1,\ldots,v_N)$ be a tuple of points in
$\mathbb R^D$. The finite integer $D$ is part of the certificate format
$i$, together with $K$ and the graph formula; it is not a quantified field
variable. Scalar triangulation in fact permits the dimension of the ambient
space containing $C_a$, but no fixed-dimension restriction is needed here.
Write
\[
 P_K(v;u)\ \Longleftrightarrow\
 \bigvee_{\sigma\in K}\exists(\lambda_j)_{j\in\sigma}
 \left(\lambda_j\ge0,\ \sum_{j\in\sigma}\lambda_j=1,
       \ u=\sum_{j\in\sigma}\lambda_jv_j\right).
\]
Let $\operatorname{Real}_K(v)$ assert affine independence on each simplex
and that two realized simplices intersect exactly in the realization of
their common face. These are finite first-order conditions; for an empty
common face the intersection condition asserts disjointness. For a fixed
formula $\varphi_i(u,z;b)$, with a fixed finite length for the additional
tuple $b$, define
\[
 \Gamma_i(a,v,b;u,z)
 \ \Longleftrightarrow\
 P_K(v;u)\ \wedge\ z\in C_a\ \wedge\varphi_i(u,z;b).
\]
Let $\operatorname{Homeo}(\Gamma_i)$ express that this relation is a
bijective graph between the indicated sets and that both it and its
inverse are continuous. Continuity is expressed in the usual
$\forall x\,\forall\varepsilon>0\,\exists\delta>0\,\forall x'$
form, using squared Euclidean distances if desired. We can now take
\[
\begin{aligned}
 &C_i(a)\ \Longleftrightarrow\\
 &\quad\exists v\,\exists b\,\exists(c_1,\ldots,c_N)\in[0,1]^N\;
 \bigl[\operatorname{Real}_K(v)\wedge\operatorname{Homeo}(\Gamma_i)\\
 &\qquad{}\wedge\!\bigwedge_{\sigma\in K}
 \forall\lambda\in\Delta_\sigma\,\forall z\\
 &\qquad\quad\bigl(\Gamma_i(a,v,b;\textstyle\sum_{j\in\sigma}\lambda_jv_j,z)\\
 &\qquad\qquad\Rightarrow
 \tau_a(z)=\textstyle\sum_{j\in\sigma}\lambda_jc_j\bigr)\bigr].
\end{aligned}
\tag{A5.17.2a}
\]
Here $\Delta_\sigma$ is the barycentric simplex defined by nonnegative
coordinates with sum one. In particular, the single-vertex cases force
$c_j=\tau_a(\pi(v_j))$. Every quantifier displayed ranges over a finite
tuple of elements of the field. The index $i$ specifies only the finite
complex, source ambient dimension, and graph formula, not the values of
their coefficients.

Shiota's exact scalar triangulation theorem~\cite[Theorem 2.2(1), p.~6]{A5-Shiota} holds in every elementary
extension of $\mathcal R$, since every such extension is again an
o-minimal expansion of a real closed field. To spell out the latter point,
for each formula in one ordered-field variable, o-minimal cell decomposition
in $\mathcal R$ supplies a uniform finite bound on the number of interval
and point components of its fibres. The sentence giving that bound, with
the other variables universally quantified, belongs to the complete theory
and hence holds in every elementary equivalent model. This applies to
every formula. Therefore each element $a$ of the parameter set in every
such extension satisfies at least one $C_i$.

We claim that finitely many $C_i$ cover $A$ already in $\mathcal R$.
If not, the type
\[
 \{a\in A\}\cup\{\neg C_i(a): i\text{ a certificate format}\}
 \tag{A5.17.3}
\]
would be finitely satisfiable in $\mathcal R$. Compactness, applied with
the elementary diagram of $\mathcal R$, would produce an elementary
extension containing a realization of this type. Scalar triangulation
there supplies a finite complex and a definable homeomorphism with a
finite parameter tuple, hence one of the listed certificate formats,
a contradiction. Thus a fixed finite list of formulas and complexes
suffices. The indexing collection is the set of all finite formulas and
finite complexes. We assume neither that this set is countable nor that
its certificates can be effectively enumerated or decided; logical
compactness applies to this arbitrary set of formulas. This argument
uses the all-model scalar theorem; it does not
infer a format bound from the pointwise real-surface approximation theorem.

\medskip\noindent\textbf{2. Finitely many abstract pairs and fixed collars.}
Use definable choice in this finite certificate family to select, for every
$a$, a certificate and its coordinates. No continuity of these choices is
asserted. Refine the finite definable partition by the set of vertices
whose chosen scalar height is zero (and, if desired, by those at height
one). There are only finitely many such patterns.

Use the standard abstract realization $|K|$ with basis-vector vertices and
the simplicial coordinate map $\rho_a:|K|\to P_a$. On each part of the
finite partition this is one definable family. Put
\[
 \Pi_a=\pi_a\circ\rho_a,
 \qquad \ell_a=\tau_a\circ\Pi_a.
\]
The domain complex is now fixed; $\ell_a$ is PL on each of its simplices
and varies only through its finitely many vertex heights. Its zero set is
one fixed subcomplex $L\subset K$ on the chosen partition part.

For every real $a$ under consideration, the classical cell-like argument
in A5.016 makes $C_a$ a topological three-manifold with its two extreme
levels as boundary. The finite link proof in A5.016 then says that $K$ is
a PL three-manifold and $L$ a union of boundary components. Consequently
each of the finitely many abstract pairs $(K,L)$ that actually occurs
admits a fixed PL collar
\[
 c_{K,L}:|L|\times[0,1]\longrightarrow|K|.
\]
Choose such a collar once for each pair. This is a finite list of choices,
not a parameter-dependent selection of arbitrarily complicated collars.
We make no assertion about cell-likeness or this topological conclusion in
an arbitrary elementary extension; only the preceding scalar-format step
used such extensions.

\medskip\noindent\textbf{3. Fixed subdivisions and rational stellar choices.}
For a fixed pair, triangulate the collar once so that its images lie in
simplices of $K$, its source height is linear, and its own graph is
simplicial. Because every $\ell_a$ is linear on the same original $K$,
this one refinement makes all $\ell_a$ linear on the required image
simplices. Its zero-vertex pattern is fixed.

Choose $\epsilon(a)>0$ below all positive source and range vertex heights
required in Lemma 1 of A5.016, and small enough that
$\ell_a^{-1}([0,\epsilon(a)])$ lies in the fixed collar image. This can
be done definably: the complement of the open collar is a fixed compact
polyhedron disjoint from $L$, and the minimum of the PL function $\ell_a$
there is a positive definable function. A common fixed triangulation
of that complement even makes this minimum the minimum of finitely many
affine combinations of the original vertex heights.

The stellar choices in the height-leveling proof can now be made with
fixed rational formulas. For a mixed simplex with zero-face barycenter
$z$ and positive-face barycenter $w$, let $d_a=\ell_a(w)>\epsilon(a)$.
Then
\[
 (1-\epsilon(a)/d_a)z+(\epsilon(a)/d_a)w
 \tag{A5.17.4}
\]
is an interior point of that simplex at the desired height. The analogous
formula uses the fixed source height on the collar domain. The same finite
list of mixed faces is starred in the same order for every $a$ on the
partition part. The resulting canonical simplicial isomorphism is a
definable family: its affine coefficients are obtained by finite linear
algebra from these finitely many moving vertices. All denominators are
nonzero because the prescribed simplices remain nondegenerate.

Thus we obtain a single definable family of height-preserving collars
\[
 C_a:|L|\times[0,\epsilon(a)]
       \longrightarrow\ell_a^{-1}([0,\epsilon(a)]).
\]
For each fixed $a$ it is continuous in both displayed variables. It need
not be continuous as $a$ varies, because the chosen scalar triangulation
$\Pi_a$ need not vary continuously.

\medskip\noindent\textbf{4. Definable choice of an error scale.}
Repeat the coordinate extraction in A5.016 for this family of collars.
The identification $G\to|L|$ at height zero is definable in $a$ because
it is the inverse of $\Pi_a$ on the zero fibre. We obtain one definable
family of homeomorphisms
\[
 h_{a,s}:F\longrightarrow G,
 \qquad 0<s\le\epsilon(a),
\]
with $h_{a,s}\to q_a$ uniformly on $F$ as $s\downarrow0$ for each fixed
$a$. The error
\[
 E(a,s)=\sup_{x\in F}\|h_{a,s}(x)-q_a(x)\|
\]
is a definable function, since $F$ is compact and the slices are
continuous. The relation
\[
 0<s\le\epsilon(a),\qquad E(a,s)<\eta
\]
has a solution for every $(a,\eta)\in A\times(0,1]$. Definable choice
provides a single function $s(a,\eta)$ satisfying it. Setting
$\psi_{a,\eta}=h_{a,s(a,\eta)}$ proves (A5.17.1)--(A5.17.2).
\end{proof}

\paragraph{What this still does not prove.}
Definable choice does not give joint continuity in $(a,\eta)$. For the
family mapping-cylinder application, the separate remaining task is to
replace $\psi$ by a jointly continuous definable family while retaining
uniform convergence to $q_a$ as the error parameter tends to zero,
locally uniformly in $a$.

A useful quantitative identity is available for sewing. If
$u,v:F\to G$ are homeomorphisms, then for $d=uv^{-1}:G\to G$,
\[
 \|d-\operatorname{id}_G\|_\infty
   =\|u-v\|_\infty
   =\|d^{-1}-\operatorname{id}_G\|_\infty.
 \tag{A5.17.5}
\]
Thus approximants close to the same possibly noninjective $q_a$ give
transition homeomorphisms close to the identity; inversion of $q_a$ is
not required. These transition maps lie in a fixed definable family
whenever $u,v$ do. The fixed-format local contraction theorem may apply
to that family, but the finite closure-pool and collar-radius construction
needed to sew all parameter cells has not been proved in this note.

A5.017 provides a uniform formula family, but not the jointly continuous
family required by the mapping-cylinder recognition criterion. It and
A5.016 are method results with no original A.5 gate movement credit.
The continuity-sewing and general-scope conclusions remain unproved here.

